\section*{Chapitre \ref{ch:b3:stem-cells} --- Cellules souches, régénération et vieillissement}

\begin{solution}{exo:b3:stem-cells:1}
\omterm{def:b3:stem-cells:stem}{Auto-renouvellement}~: une division qui produit au moins une fille encore
souche --- les cellules \emph{Lgr5} de la base de la crypte. Potence~:
l'éventail des destins qu'une cellule peut encore adopter --- la \omterm{def:b3:stem-cells:stem}{cellule
souche} de la crypte est multipotente et fait des cellules absorbantes,
caliciformes, entéroendocrines, de Paneth et en touffe. \omterm{def:b3:stem-cells:stem}{Niche}~: les
\omterm{prop:b3:stem-cells:crypt}{cellules de Paneth} et leurs signaux Wnt, Notch et EGF, qui tiennent le
caractère souche en place. \omterm{def:b3:stem-cells:stem}{Cellule d'amplification transitoire}~: une
fille engagée qui se divise encore quatre ou cinq fois dans la crypte
avant de se différencier en remontant la villosité.
\end{solution}

\begin{solution}{exo:b3:stem-cells:2}
Des cellules de moelle isolées ont produit des colonies spléniques
contenant globules rouges, granulocytes et plaquettes~: une cellule,
plusieurs lignées --- la multipotence, et la réponse linéaire à la dose
a montré que chaque colonie venait d'une seule cellule.
L'\omterm{def:b3:stem-cells:stem}{auto-renouvellement} a exigé une seconde souris~: des cellules d'une
colonie, injectées à un autre receveur irradié, ont fait de nouvelles
colonies, donc la cellule fondatrice avait produit des filles douées de
sa propre capacité.
\end{solution}

\begin{solution}{exo:b3:stem-cells:3}
\emph{Oct4}, \emph{Sox2}, \emph{Klf4}, \emph{c-Myc}. La \omterm{def:b3:chromatin-epigenetics:nucleosome}{chromatine} du
fibroblaste a refermé les gènes de la \omterm{def:b3:stem-cells:pluripotency}{pluripotence} derrière
l'\omterm{def:b3:chromatin-epigenetics:nucleosome}{hétérochromatine} et la \omterm{def:b3:chromatin-epigenetics:methylation}{méthylation de l'ADN}, de sorte que les facteurs
trouvent d'abord peu de sites accessibles et doivent recruter des
enzymes de remodelage au fil de nombreux \omterm{def:b3:cell-cycle-apoptosis:cdk}{cycles cellulaires}~; la cellule
doit aussi éteindre son propre programme, échapper à la \omterm{def:b3:cell-cycle-apoptosis:other}{sénescence} et
traverser un état intermédiaire stochastique où la plupart des cellules
calent ou meurent. Seule la rare cellule où chaque étape réussit émerge
pluripotente, des semaines plus tard.
\end{solution}

\begin{solution}{exo:b3:stem-cells:4}
La planaire garde des \omterm{def:b3:stem-cells:stem}{cellules souches} adultes pluripotentes, les
\omterm{def:b3:stem-cells:regeneration}{néoblastes}, qui migrent vers la plaie et bâtissent ce qui manque. La
salamandre n'a pas un tel réservoir~: des cellules différenciées du
moignon se dédifférencient en \omterm{def:b3:stem-cells:stem}{progéniteurs} qui, avec les \omterm{def:b3:stem-cells:stem}{cellules
souches} résidentes du tissu, forment un \omterm{def:b3:stem-cells:regeneration}{blastème} qui réexécute le
programme du membre. Les greffes de tissus marqués un à un par Kragl ont
montré que les cellules du \omterm{def:b3:stem-cells:regeneration}{blastème} restent fidèles à leur origine ---
le muscle fait du muscle, le cartilage du cartilage --- de sorte que le
\omterm{def:b3:stem-cells:regeneration}{blastème} est un mélange de \omterm{def:b3:stem-cells:stem}{progéniteurs} restreints à leur lignée, non
une masse pluripotente.
\end{solution}

\begin{solution}{exo:b3:stem-cells:5}
$(11 - 5)/0.075 = 80$ doublements. En partant de \qty{8}{kb}~:
$(8 - 5)/0.075 = 40$. Avec \qty{60}{\%} de compensation la perte nette
vaut \qty{30}{bp}~: $6000/30
= 200$.
\end{solution}

\begin{solution}{exo:b3:stem-cells:6}
Quatorze divisions de \omterm{def:b3:stem-cells:stem}{cellules souches} par jour donnent 28 filles~; 14
doivent rester souches, donc 14 s'engagent --- une fille engagée par
\omterm{def:b3:stem-cells:stem}{cellule souche} et par jour. Chacune donne $2^{4} = 16$ cellules~:
$14\times 16 = 224$ cellules de villosité par jour, l'ordre des
\num{300} observées (une cinquième division de transit en donne
\num{448}). L'étage souche fournit l'équivalent de quatorze cellules en
divisions, l'étage de transit les \num{200} autres.
\end{solution}

\begin{solution}{exo:b3:stem-cells:7}
$2\times 10^{11}/2^{20} = 1.9\times 10^{5}$ filles engagées par jour~;
avec \num{10000} \omterm{def:b3:stem-cells:stem}{cellules souches}, 19 par \omterm{def:b3:stem-cells:stem}{cellule souche} et par jour ---
contre une division par an. Le compte de vingt divisions vaut par lignée
de la \omterm{def:b3:stem-cells:stem}{cellule souche} au globule rouge, mais les \omterm{def:b3:stem-cells:stem}{progéniteurs}
intermédiaires ne servent pas une seule fois~: les réservoirs de
\omterm{def:b3:stem-cells:stem}{progéniteurs} multipotents et engagés s'entretiennent des semaines durant
par leurs propres divisions, de sorte que presque toute la division se
passe aux étages de \omterm{def:b3:stem-cells:stem}{progéniteurs} et qu'on fait rarement appel à la
\omterm{def:b3:stem-cells:stem}{cellule souche}.
\end{solution}

\begin{solution}{exo:b3:stem-cells:8}
$\mu(50) = 10^{-3}\mathrm{e}^{1.74} = 5.7\times 10^{-3}$~; $\mu(70) =
10^{-3}\mathrm{e}^{3.48} = 0.032$~; $\mu(90) = 10^{-3}\mathrm{e}^{5.22} =
0.19$ par an. Avec $\mu_{0}/\gamma = 0.0115$~: $S(70) = \exp[-0.0115
\times 31.5] = 0.70$, $S(90) = \exp[-0.0115\times 184] = 0.12$. Médiane~:
$30 + \ln(61)/0.087 = 77$. En divisant $\mu_{0}$ par deux~:
$\ln(121)/0.087 = 55$, médiane 85 --- un temps de doublement, huit ans,
de gagné. En divisant $\gamma$ par deux, à $0.0435$~: $\ln(31)/0.0435 = 79$,
médiane 109~: ralentir l'exposant rapporte quatre fois plus que diviser
la constante.
\end{solution}

\begin{solution}{exo:b3:stem-cells:9}
Le tiers restant doit tripler~: $\log_{2}3 = 1.6$ division par cellule.
\omterm{def:b3:stem-cells:regeneration}{Croissance compensatrice}, parce que les lobes existants grossissent par
division de cellules sur place~; les lobes retirés, avec leurs canaux,
leurs vaisseaux et leur architecture, ne sont pas rebâtis --- le foie
retrouve sa masse et sa fonction, non sa forme.
\end{solution}

\begin{solution}{exo:b3:stem-cells:10}
À chaque événement l'effectif $k$ du clone passe à $k + 1$ ou $k - 1$
avec la même probabilité, de sorte que son effectif attendu ne change
jamais~; quand la marche s'achève il vaut $N$ (prise de contrôle) ou $0$
(perte), et l'espérance $N\,P + 0 = k$ donne $P = k/N$. Une seule
cellule~: $1/14$. Une mutation neutre dans une \omterm{def:b3:stem-cells:stem}{cellule souche} se fixe
dans sa crypte avec la probabilité $1/14$ et se perd sinon en quelques
mois --- la dérive élimine treize mutations sur quatorze.
\end{solution}

\begin{solution}{exo:b3:stem-cells:11}
La marche est maintenant biaisée~: le clone gagne une place plus souvent
qu'il n'en perd, de sorte que son effectif attendu croît et que la
probabilité de prise de contrôle monte vers la certitude à mesure que
$k$ grandit --- pour un avantage de \qty{10}{\%} dans une \omterm{def:b3:stem-cells:stem}{niche} de
quatorze, à peu près le triple de $1/14$ pour une seule cellule, et bien
plus une fois que le clone tient quelques places. Au fil des décennies,
les \omterm{def:b3:stem-cells:stem}{cellules souches} porteuses de telles mutations (\emph{DNMT3A},
\emph{TET2} dans la moelle) prennent leurs \omterm{def:b3:stem-cells:stem}{niches}, et à soixante-dix ans
une personne sur cinq tire une large part de son sang d'un seul clone.
Ce n'est pas un \omterm{def:b3:cancer-biology:hallmarks}{cancer}~: les cellules se différencient encore et
obéissent à leur hiérarchie. Mais la mutation suivante dispose désormais
d'un clone de milliards de cellules où apparaître, et le risque de
leucémie est décuplé.
\end{solution}

\begin{solution}{exo:b3:stem-cells:12}
Médiane $= 30 + \gamma^{-1}\ln(1 + \gamma\ln 2/\mu_{0})$~: pour $\mu_{0} =
10^{-2}$, $30 + \ln 7/0.087 = 52$~; $10^{-3}$~: 77~; $10^{-4}$~: $30 +
\ln 604/0.087 = 104$. Dans la forme de Gompertz pure chaque division par
dix ajoute les mêmes $\ln 10/\gamma = 26$ ans. Les rendements
décroissent en pratique parce que les gains historiques sont venus de la
suppression d'un terme indépendant de l'âge --- infection, accouchement,
accident, la constante $A$ de Makeham dans $\mu = A
+ \mu_{0}\mathrm{e}^{\gamma t}$ --- et qu'une fois $A$ petit devant
l'exponentielle aux âges qui comptent, le réduire encore ne change
presque rien~; le $\mu_{0}$ intrinsèque a beaucoup moins bougé. Une
baisse de \qty{20}{\%} de $\gamma$ (à $0.070$) donne
$30 + \ln 49/0.070 = 86$~: neuf ans, gagnés à tout âge plutôt qu'en
repoussant quelques causes de mort.
\end{solution}

\begin{solution}{pb:b3:stem-cells:1}
\textbf{1.} $3.5\times 10^{11}$ par jour, $4\times 10^{6}$ par seconde.
\textbf{2.} $2^{20} \approx 10^{6}$ cellules par fille engagée~;
$3.5\times 10^{11}/10^{6} = 3.3\times 10^{5}$ filles engagées par jour.
\textbf{3.} $33$ par \omterm{def:b3:stem-cells:stem}{cellule souche} et par jour, \num{12000} par an.
\textbf{4.} Les réservoirs de \omterm{def:b3:stem-cells:stem}{progéniteurs} doivent s'entretenir~: les
\omterm{def:b3:stem-cells:stem}{progéniteurs} multipotents et engagés se divisent des semaines durant et
maintiennent leur propre effectif, de sorte que presque toutes les vingt
divisions se passent à des étages qui se renouvellent eux-mêmes, et que
la \omterm{def:b3:stem-cells:stem}{cellule souche} ne fournit une fille engagée que rarement --- de
l'ordre d'une fois par an.
\textbf{5.} $10 - 20\times 0.06 = \qty{8.8}{kb}$~: bien au-dessus de
$L_{\text{crit}}$, et le globule rouge ne se divise plus jamais~; la
réserve ne compte que pour les \omterm{def:b3:stem-cells:stem}{progéniteurs} qui resservent.
\textbf{6.} Perte nette $0.3\times 60 = \qty{18}{bp}$~; $6000/18 = 333$
divisions~; à une par an, $333$ ans --- bien au-delà d'une vie, de sorte
que les télomères ne limitent pas la \omterm{def:b3:stem-cells:stem}{cellule souche} elle-même~; d'autres
dégâts le font.
\textbf{7.} $14$ divisions par jour~; $7$ remplacent des \omterm{def:b3:stem-cells:stem}{cellules souches}
perdues et $7$ s'engagent~; $7\times 2^{4} = 112$ cellules de villosité
par jour.
\textbf{8.} $N^{2} = 196$ événements par cellule à raison d'un par jour~:
environ \qty{200}{d}, soit six à sept mois --- le temps sur lequel
les cryptes à confetti sont devenues monochromes.
\textbf{9.} $1/14 \approx \qty{7}{\%}$~; $10^{7}/14 \approx 7\times 10^{5}$
cryptes.
\textbf{10.} La dérive rejette treize mutations sur quatorze en quelques
mois. Une \omterm{def:b3:stem-cells:stem}{division asymétrique} à vie garderait chaque \omterm{def:b3:stem-cells:stem}{cellule souche}
mutée pour toujours, de sorte que les mutations s'accumuleraient dans
chaque lignée sans jamais être purgées.
\textbf{11.} Le caractère souche est une position dans la \omterm{def:b3:stem-cells:stem}{niche} et un
jeu de signaux, non une propriété permanente d'une cellule~: toute
cellule qui atteint les signaux de Paneth peut en tenir le rôle.
\textbf{12.} Induire un marquage multicolore dans les \omterm{def:b3:stem-cells:stem}{cellules souches} à
un instant et suivre les cryptes. Dérive symétrique~: les cryptes
deviennent monochromes en quelques mois. \omterm{def:b3:stem-cells:stem}{Division asymétrique}~: chaque
crypte garde indéfiniment toutes ses couleurs, chacune en fin ruban le
long de la villosité.
\textbf{13.} $6000/60 = 100$ doublements.
\textbf{14.} Environ $75$ restants~: la cellule compte des divisions et
non le temps du calendrier, puisque la décennie au congélateur n'a rien
coûté.
\textbf{15.} $(8000 - 4000)/30 = 133$ ans après vingt ans --- à 153 ans.
La longueur moyenne des télomères des leucocytes ne limite pas
elle-même la vie~; les télomères les plus courts, et les autres
marqueurs, comptent davantage.
\textbf{16.} La \omterm{def:b3:cancer-biology:telomerase}{télomérase} lève une barrière, la \omterm{thm:b3:stem-cells:hayflick}{sénescence
réplicative}, et se retrouve donc dans presque tous les \omterm{def:b3:cancer-biology:hallmarks}{cancers}, qui ont
besoin d'une division illimitée~; mais les fibroblastes dotés de
\omterm{def:b3:cancer-biology:telomerase}{télomérase} ont gardé leurs points de contrôle, leur inhibition de
contact et un p53 intact, de sorte que la seule division illimitée n'en
a pas fait des \omterm{def:b3:cancer-biology:hallmarks}{cancers}. Nécessaire, non suffisante.
\textbf{17.} $\mu(40) = 2.4\times 10^{-3}$, $\mu(60) = 0.014$, $\mu(80) =
0.077$, $\mu(100) = 0.44$ par an. $S = \exp[-0.0115(\mathrm{e}^{\gamma
(t-30)} - 1)]$~: $0.98$ à 40 ans, $0.87$ à 60, $0.42$ à 80 et $0.006$ à
100.
\textbf{18.} Médiane $30 + \ln 61/0.087 = 77$. Dix pour cent survivent
quand $\mathrm{e}^{\gamma(t-30)} = 1 + \gamma\ln 10/\mu_{0} = 201$~: $t = 30 +
5.3/0.087 = 91$.
\textbf{19.} \qty{0.67}{mm/d}~; $\log_{2}1000 = 10$ doublements en 60
jours, un tous les six jours.
\textbf{20.} Environ $10^{8}$ divisions, donc $0.1$ mutation \omterm{def:b3:cancer-biology:genes}{oncogène} par
\omterm{def:b3:stem-cells:regeneration}{régénération} --- une sur dix membres. Les axolotls ne font presque jamais
de \omterm{def:b3:cancer-biology:hallmarks}{tumeurs}~: à sang froid et lents, avec une répression tumorale
robuste, et les cellules du \omterm{def:b3:stem-cells:regeneration}{blastème} restent restreintes à leur lignée
et se redifférencient au lieu de continuer à se diviser~; un tissu en
\omterm{def:b3:stem-cells:regeneration}{régénération} étouffe même les \omterm{def:b3:cancer-biology:hallmarks}{tumeurs} induites. Un mammifère a plus de
cellules, une température et un métabolisme plus élevés, et une vie plus
longue au cours de laquelle un clone échappé pourrait grandir.
\textbf{21.} Dénerver le moignon et y implanter une bille libérant la
protéine candidate, ou greffer un tissu conditionné par un nerf~: une
\omterm{def:b3:stem-cells:regeneration}{régénération} rétablie signale un facteur diffusible. Candidats
identifiés chez le triton et l'axolotl~: la protéine nAG du gradient
antérieur, la neuréguline-1 et les FGF.
\textbf{22.} La \omterm{prop:b3:stem-cells:crypt}{signalisation Wnt} est élevée à la queue et faible à la
tête, et la plaie de l'extrémité antérieure d'un morceau fait
normalement une tête parce que Wnt y est faible~; invalider un
inhibiteur de Wnt élève Wnt partout, les deux plaies lisent «~postérieur~»,
et deux queues poussent. Le morceau connaît son propre axe grâce au
gradient résiduel de son muscle, et la plaie rétablit le gradient par
rapport à lui.
\textbf{23.} Le tiers restant doit tripler~: chaque hépatocyte se divise
une fois (pour atteindre les deux tiers) et la moitié se divise à
nouveau --- une moyenne de $1.6$ division. Les lobes survivants enflent
jusqu'à la masse initiale~; les lobes manquants, leurs canaux et leurs
vaisseaux ne sont pas rebâtis.
\textbf{24.} Une \omterm{def:b3:stem-cells:regeneration}{régénération} qui remplacerait les tissus usés et
abîmés ralentirait l'accumulation que mesure l'exposant, ce qui
abaisserait $\gamma$ et aplatirait la courbe --- les axolotls montrent
peu de \omterm{def:b3:cell-cycle-apoptosis:other}{sénescence}. Le prix serait un terme constant plus élevé, dû aux
\omterm{def:b3:cancer-biology:hallmarks}{cancers} dans un corps chaud, grand et de longue vie.
\textbf{25.} Compte de Hayflick d'environ $100$ doublements~; une \omterm{def:b3:stem-cells:stem}{cellule
souche} aidée par la \omterm{def:b3:cancer-biology:telomerase}{télomérase} dispose de quelque $330$ divisions, plus
de trois siècles à une par an~; espérance de vie médiane de $77$ ans.
\end{solution}
